Capable large language models --- GPT-4, Claude 3, DeepSeek-V3, Qwen3, LLaMA-3 --- have rekindled a long-standing research program: building societies of autonomous agents that solve problems through interaction. An \textbf{LLM-based multi-agent system (LLM-MAS)} is a software system in which two or more agents, each whose policy is parameterized in significant part by an LLM, exchange natural-language messages over a directed communication graph to accomplish a task or simulate a social process. The formulation was popularized by five canonical 2023 systems: \textbf{CAMEL} (Li et al., 2023), \textbf{AutoGen} (Wu et al., 2023), \textbf{MetaGPT} (Hong et al., 2023), \textbf{ChatDev} (Qian et al., 2023), and Stanford's \textbf{Generative Agents} (Park et al., UIST 2023). The area has since matured from exploratory curiosity into a mainstream research field, with at least seven major surveys between 2024 and 2026 (Wang et al., 2024 \emph{FCS}; Guo et al., 2024; Cheng et al., 2024; Li et al., 2024 \emph{Vicinagearth}; Chen et al., 2024; Aratchige \& Ilmini, 2025; Yan et al., 2025). This survey consolidates the conceptual landscape, algorithmic primitives, dominant frameworks and applications, evaluation protocols, and open problems of LLM-MAS into a single retrieval-friendly reference. The \textbf{scope} is bounded by Section 1.3: we cover language-native message-passing societies of two or more LLM agents, excluding single-agent persona prompting and pure MARL. The \textbf{taxonomy} of Section 3 organizes the field along three or\ldots{}
